\newpage{\ } 
\thispagestyle{empty}
\chapter*{Resumen}

El presente Trabajo Fin de Grado tiene como objetivo analizar, diseñar y desarrollar una aplicación web progresiva orientada a centralizar la actividad comercial de un distribuidor de productos capilares profesionales. El caso de estudio parte de una necesidad concreta: gestionar de forma continua clientes, visitas comerciales, pedidos, entregas, cobros, catálogo de productos e información administrativa asociada. Actualmente, estos procesos se apoyan en herramientas independientes, como hojas de cálculo, bases de datos locales y programas de facturación. Esta situación provoca una fragmentación de la información, incrementa la dependencia de tareas manuales y dificulta mantener una trazabilidad clara de la actividad diaria.

Ante este contexto, la solución propuesta buscará ordenar la gestión comercial y servir como base para la digitalización progresiva de la relación entre el distribuidor, los comerciales y las peluquerías profesionales. Para ello, se estudiará el funcionamiento actual del negocio, se identificarán sus principales limitaciones y se definirán los requisitos funcionales, no funcionales y de información necesarios para plantear una solución modular y extensible.

La aplicación se planteará como un sistema web full-stack basado en Next.js, React, TypeScript y PostgreSQL, organizado en torno a distintos perfiles de usuario: administrador, comerciales y clientes profesionales. Esta diferenciación permitirá adaptar las funcionalidades del sistema a las necesidades de cada rol, incorporando autenticación, trazabilidad de accesos, control de permisos, gestión de usuarios y perfil, administración de clientes y visitas comerciales, planificación operativa con estimación de rutas, consulta y mantenimiento del catálogo de productos y coloración, así como un flujo completo de pedidos que abarcará la preparación de repartos, la entrega por comercial o agencia, la validación mediante QR y la gestión de pagos parciales. Además, la aplicación contemplará herramientas de seguimiento técnico interno del salón, comunicaciones comerciales, auditoría exportable, configuración de avisos automáticos, gestión del cargo de agencia y una base técnica preparada para rate limiting, soporte PWA e integración futura con Factusol mediante n8n.

Con este trabajo se pretende proporcionar una base sólida para reducir la dependencia de procesos manuales, mejorar la organización de la información y facilitar la evolución futura de la plataforma hacia una solución integral de apoyo a la actividad comercial del distribuidor.

\vspace{.5cm}

\noindent \textbf{Palabras clave:} aplicación web progresiva, gestión comercial, distribución B2B, productos capilares, digitalización de procesos, gestión de clientes, gestión de pedidos, sistemas de información, autenticación, autorización, roles de usuario, trazabilidad, arquitectura web, desarrollo full-stack, Next.js, PostgreSQL.